% Generated draft: requires mathematical and source review.
\documentclass[11pt]{article}
\usepackage[margin=25mm]{geometry}
\usepackage{fontspec}
\setmainfont{Latin Modern Roman}
\usepackage{amsmath,amssymb,mathtools,unicode-math}
\setmathfont{Latin Modern Math}
\usepackage{enumitem,needspace,xurl,hyperref}
\setlength{\parindent}{0pt}
\setlength{\parskip}{5pt}
\setlength{\emergencystretch}{3em}
\newcommand{\N}{\mathbb{N}}
\newcommand{\Z}{\mathbb{Z}}
\newcommand{\Q}{\mathbb{Q}}
\newcommand{\R}{\mathbb{R}}
\newcommand{\C}{\mathbb{C}}
\newcommand{\F}{\mathbb{F}}
\newcommand{\A}{\mathbb{A}}
\newcommand{\PP}{\mathbb{P}}
\newcommand{\E}{\mathbb{E}}
\newcommand{\eps}{\varepsilon}
\newcommand{\dd}{\,\mathrm{d}}
\begin{document}
\section*{UnsolvedMath ID: 30005015}
\textbf{Source problem code:} catalogue entry
\subsection{Definitions and mathematical statement}
The supplied record asks for the dimension of a planar spline space, with notation apparently intended to be \(S_m^r(\Delta)\). The excerpt begins after the definition of this space, and its typography separates the superscript \(r\) from the remaining notation. Consequently, its coefficient field, degree convention, smoothness convention, admissible subdivisions, and boundary conditions cannot be recovered completely from the supplied text.

For an explicitly conditional formulation, suppose the intended convention is the following: \(\Delta\) is a finite planar polygonal subdivision of a domain \(\Omega\subseteq\R^2\), \(m,r\in\Z_{\geq 0}\), and
\[
 S_m^r(\Delta)=\left\{f\in C^r(\Omega;\R):\ \text{for every two-dimensional cell }\sigma\in\Delta,\ \exists p_\sigma\in\R[x,y],\ \deg p_\sigma\leq m,\ f|_\sigma=p_\sigma|_\sigma\right\}.
\]
Here degree means total degree, and the displayed convention imposes no additional boundary conditions. These are provisional conventions for explaining the target, not hypotheses established by the supplied excerpt; the meaning of smoothness at the boundary must also be specified if this model is adopted.

\textbf{Requested characterization.} For the intended class of subdivisions \(\Delta\) and intended degree and smoothness parameters \(m,r\), determine \(\dim S_m^r(\Delta)\) in terms of the combinatorics of \(\Delta\), its local geometry, and its global geometry. Under the conditional real-valued convention above, this means determining \(\dim_{\R}S_m^r(\Delta)\). The source does not specify which geometric data are allowed, what form the characterization should take, or the ranges of parameters for which it is requested. Thus the supplied record supports this characterization problem but does not determine a unique fully quantified formulation.

The original excerpt subsequently changes topic to an interpolation question attributed to Shekhtman. That appended question is distinct from the spline-dimension target and supplies no missing spline hypotheses.
\subsection{Short English statement}
Determine the dimension of the intended planar spline space from the subdivision's combinatorics and its local and global geometry. The supplied excerpt omits the space's defining hypotheses, so the displayed conventional model is conditional and requires human review. The record's open status and verification labels are supplied curation claims.
\subsection{Sources}
\begin{itemize}
\item Supplied local UnsolvedMath registry record, permanent ID \(30005015\), title: Geometry and Combinatorics of Planar Spline Dimensions; exact source problem code: \texttt{catalogue entry}. No proposer is identified for the spline question.
\item Supplied source citation: \emph{Mini-Workshop: Interpolation, Approximation, and Algebra (2023)}, Oberwolfach Reports, pp.~383--428, DOI \texttt{10.4171/owr/2022/7}. \url{https://doi.org/10.4171/owr/2022/7}. The record describes Oberwolfach Reports as documenting mathematical workshops and their open-problem discussions.
\item Supplied literature source: Schenck--Stillman--Yuan, spline bounds. \url{https://arxiv.org/abs/1909.13399}. No more detailed bibliographic identification or applicable theorem is supplied.
\item Supplied literature source: OWR spline report. \url{https://doi.org/10.4171/owr/2022/7}.
\item Literature assessment dated exactly \texttt{2026-08-22}: the supplied curation labels the problem open and states that known formulas and homological bounds do not provide the requested combinatorial-geometric characterization in full generality. This is supplied dated curation derived from the OWR record, not an independently verified assessment; the record expressly requests independent verification of time-sensitive status and cited claims.
\item The registry labels the statement \texttt{corrected\_verified} and reports that it was checked against the official Oberwolfach report. These are supplied verification claims; the report and cited literature have not been independently consulted for this draft.
\end{itemize}
\end{document}
